%% d2b_classification.tex -- DISTILLED notes: classification models
%% (logistic regression, LDA, QDA, naive Bayes) -- Slide Set 3 pp. 1-33 & 43-93,
%% Lab 3, HW 3 (Q4-Q12), practice exam Q14/Q15/Q16/Q20.
\section{Classification models: logistic regression, LDA, QDA, naive Bayes}
\ixh{classification}

%% ==================================================================
\subsection{Classification set-up; why not linear regression?}
\ixh{classification}

\begin{topic}{Classification in general}
\begin{itemize}
\item Predicting a \ix{qualitative response} for an observation $=$ \ix{classification}: assign the
  observation to a class. Usually we first predict the \emph{probability} of each class
  (a quantitative number) and then classify \src{S3.4}.
\item Classifiers in this course: \ix{KNN}, \ix{logistic regression}, \ix{discriminant analysis}
  (\ix{LDA}, \ix{QDA}), \ix{naive Bayes}.
\item Running example: \ixc{Default} data (ISLR): $Y=$ \texttt{default} (Yes/No);
  $X_1=$ annual \texttt{income}, $X_2=$ monthly credit-card \texttt{balance}, $X_3=$ \texttt{student}.
  Defaulters have much higher balances; income barely differs \src{S3.8}.
\end{itemize}
\end{topic}

\begin{keybox}[Why not linear regression for a categorical Y?]
\begin{itemize}
\item \textbf{Binary} $Y\in\{0,1\}$: \ix{linear regression} models
  $\Pr(Y=1\mid X)=\beta_0+\beta_1X_1+\dots+\beta_pX_p$, an unbounded line
  $\Rightarrow$ fitted ``probabilities'' can be $<0$ (Default data: \textbf{negative probability} for
  balance $\lesssim 500$) or $>1$. Logistic regression's S-curve stays in $[0,1]$ \src{S3.5}.
\item \textbf{$>2$ classes}: coding $Y=1$ (stroke), 2 (drug overdose), 3 (seizure) imposes an
  \textbf{ordering} and equal spacing (stroke$\to$overdose gap $=$ overdose$\to$seizure gap); a different
  coding (1 stroke, 2 seizure, 3 overdose) gives a different model; $\{$dog, cat, fish$\}$ has no
  order at all \src{S3.6}.
\item Course focuses on \textbf{binary} classification (multi-class logistic regression exists) \src{S3.7}.
\end{itemize}
\end{keybox}

%% ==================================================================
\subsection{The logistic regression model}
\ixh{logistic regression}

\begin{keybox}[Logistic regression: three equivalent forms]
With $p(X)=\Pr(Y=1\mid X)$ \src{S3.10--11, S3.24}:
\begin{align*}
p(X)&=\frac{e^{\beta_0+\beta_1X}}{1+e^{\beta_0+\beta_1X}}\in[0,1] &&\text{\ix{logistic function}}\\[2pt]
\frac{p(X)}{1-p(X)}&=e^{\beta_0+\beta_1X}\in[0,\infty) &&\text{\ix{odds}}\\[2pt]
\log\Bigl(\frac{p(X)}{1-p(X)}\Bigr)&=\beta_0+\beta_1X\in(-\infty,\infty) &&\text{\ix{log odds} $=$ \ix{logit}}
\end{align*}
\textbf{The logit is linear in $X$} (not $p$!). Multiple predictors: replace $\beta_0+\beta_1X$ by
$\beta_0+\beta_1X_1+\dots+\beta_pX_p$ everywhere \src{S3.24}. $\log=\ln$ (natural log).
Equivalent: $p=\frac{1}{1+e^{-z}}$ with $z=\beta_0+\beta_1X$ (the \ix{sigmoid}).

\textbf{Odds} $=p/(1-p)$, \emph{not} $p$: rolling a 3 on a fair die has $p=1/6$, odds $=\frac{1/6}{5/6}=1/5$
\src{S3.10}. $p=0.8\Rightarrow$ odds $=4$ (``4 to 1'').
\end{keybox}

\noindent
\begin{minipage}[c]{0.56\linewidth}
\begin{tikzpicture}
\begin{axis}[scale only axis,width=0.80\linewidth,height=2.3cm,
  xmin=-6,xmax=6,ymin=0,ymax=1,axis lines=left,
  xtick={-6,-3,0,3,6},ytick={0,0.5,1},
  xlabel={$z=\beta_0+\beta_1X$ (logit)},ylabel={$p$},
  ylabel style={yshift=-6pt},xlabel style={yshift=3pt},
  samples=120,domain=-6:6]
\addplot[very thick,dblue] {1/(1+exp(-x))};
\addplot[densely dashed,gray] coordinates {(-6,0.5) (0,0.5) (0,0)};
\node[font=\tiny,anchor=west,align=left] at (axis cs:-5.8,0.83) {$p=\frac{e^z}{1+e^z}$};
\node[font=\tiny,anchor=west] at (axis cs:0.3,0.12) {$z=0\Leftrightarrow p=0.5$};
\end{axis}
\end{tikzpicture}
\end{minipage}\hfill
\begin{minipage}[c]{0.42\linewidth}\footnotesize
\begin{tabular}{@{}rrr@{}}
\toprule
$p$ & odds & log odds\\ \midrule
0 & 0 & $-\infty$\\
0.269 & 0.368 & $-1$\\
0.5 & 1 & 0\\
0.731 & 2.72 & 1\\
0.88 & 7.33 & 2\\
0.95 & 19 & 3\\
1 & $\infty$ & $+\infty$\\ \bottomrule
\end{tabular}\\[1pt]
{\scriptsize $p$ on $[0,1]$ $\mapsto$ log odds on $(-\infty,\infty)$ \src{S3.16--21}}
\end{minipage}

\begin{topic}{Interpreting logistic coefficients}
\src{S3.12--13, S3.25}
\begin{itemize}
\item $\beta_1$ $=$ the \textbf{additive} change in the \ix{logit} (log odds) for a 1-unit increase in $X$.
\item $e^{\beta_1}$ $=$ the \textbf{multiplicative} change in the \ix{odds}:
  $\text{odds}(X{+}1)/\text{odds}(X)=e^{\beta_1}$ (the \ix{odds ratio}). A $c$-unit increase multiplies the odds by
  $e^{c\beta_1}$ (HW3 Q12: $+2$ units $\Rightarrow\times e^{2\beta_1}$). \% change in odds $=100(e^{\beta_1}-1)\%$.
\item $\beta_1$ is \textbf{NOT} the change in $\Pr(Y=1\mid X)$ per unit of $X$. Only the \textbf{sign} carries over:
  $\beta_1>0\Rightarrow X\nearrow\Rightarrow p\nearrow$;\; $\beta_1<0\Rightarrow X\nearrow\Rightarrow p\searrow$;\;
  $\beta_1=0\Rightarrow$ no association.
\item The rate of change of $p$ per unit $X$ \textbf{depends on the current value of $X$} --- not linear
  (calculus: $dp/dX=\beta_1p(1-p)$, steepest, $=\beta_1/4$, at $p=0.5$).
\item Multiple logistic regression: $\beta_j$ $=$ change in logit per unit $X_j$
  \textbf{conditional on (holding fixed) the other predictors} \src{S3.25}.
\item $\beta_0$ $=$ log odds when all $X$'s $=0$. Regression on observational data shows
  \textbf{association, not causation}\ixh{causation}.
\end{itemize}
\end{topic}

\begin{trap}[PE Q16: coefficient 0.04 is NOT ``+4\% probability'']
\ixh{coefficient interpretation}Logistic regression for cancelling a subscription; coefficient of \texttt{MonthlyFee} $=0.04$.\\
\textbf{Wrong:} ``A \$1 increase in the monthly fee \emph{causes} a customer's probability of cancelling to
increase by 4\%.'' Problems: (1) it treats the probability as linear in the fee --- the relationship is
linear in the \ix{log odds}, and the change in $p$ depends on the current fee; (2) ``causes'' ---
only \textbf{association} can be claimed.\\
\textbf{Right:} ``A \$1 increase in the monthly fee is associated with a 0.04 increase in the log odds of
cancelling'' (equivalently, odds $\times e^{0.04}\approx1.041$, about $+4.1\%$ in the \emph{odds}), or
``a higher fee is associated with a higher probability of cancelling'' (sign is positive) \src{PE Q16}.
\end{trap}

\begin{keybox}[HW3 Q11--Q12 (231): logit is the inverse of the logistic]
$p=\frac{e^z}{1+e^z}\Rightarrow p+pe^z=e^z\Rightarrow p=e^z(1-p)\Rightarrow e^z=\frac{p}{1-p}
\Rightarrow z=\ln\frac{p}{1-p}$ \src{HW3 Q11}.\quad
$z=\beta_0+\beta_1x_1$: odds at $x_1{+}2$ $=e^{\beta_0+\beta_1(x_1+2)}=e^{\beta_0+\beta_1x_1}\cdot e^{2\beta_1}$, so the odds are
multiplied by $e^{2\beta_1}$. If $\beta_1<0$: $x_1\to\infty\Rightarrow z\to-\infty\Rightarrow p\to0$;\;
$x_1\to-\infty\Rightarrow z\to+\infty\Rightarrow p=\frac{1}{1+e^{-z}}\to1$ \src{HW3 Q12}.
\end{keybox}

%% ==================================================================
\subsection{Estimating the coefficients: maximum likelihood}
\ixh{maximum likelihood}

\begin{topic}{Fitting with glm()}
\begin{itemize}
\item Coefficients are estimated by \ix{maximum likelihood} (MLE), \emph{not} least squares \src{S3.26}:
  choose $\hat\beta$'s so $\hat p(x_i)$ is near 1 when $y_i=1$ and near 0 when $y_i=0$, i.e.\ maximise
  $\ell(\beta_0,\beta_1)=\prod_{i:y_i=1}p(x_i)\prod_{i':y_{i'}=0}\bigl(1-p(x_{i'})\bigr)$ (ISLR 4.3.2).
\item Base R: \ixc{glm()} with \ixc{family = binomial} (``generalized linear model''):
\end{itemize}
\begin{rcode}
glm.fit <- glm(default ~ student + balance + income,      (*@\ixch{glm()}\ixch{Default}@*)
               data = Default, family = binomial)        (*@\ixch{family = binomial}@*)
summary(glm.fit)                                          (*@\ixch{summary()}@*)
\end{rcode}
\begin{routput}
              Estimate Std. Error z value Pr(>|z|)
(Intercept) -1.087e+01  4.923e-01 -22.080  < 2e-16 ***
studentYes  -6.468e-01  2.363e-01  -2.738  0.00619 **
balance      5.737e-03  2.319e-04  24.738  < 2e-16 ***
income       3.033e-06  8.203e-06   0.370  0.71152
    Null deviance: 2920.6  on 9999  degrees of freedom
Residual deviance: 1571.5  on 9996  degrees of freedom
AIC: 1579.5
\end{routput}
\begin{itemize}
\item \texttt{z value} $=$ Estimate$/$Std.\ Error (a \ix{z-statistic}, not $t$); small \ix{p-value}
  $\Rightarrow$ reject $H_0{:}\ \beta_j=0$ given the other predictors. Here income is not significant
  ($p=0.71$) once balance and student are in the model.
\item \texttt{studentYes} is the \ix{dummy variable} for student $=$ Yes (\ix{reference level} $=$ No).
\item Units: \texttt{income} coefficient is $3.033\times10^{-6}$ per \$1; the slide tables show
  $0.0030$ because there income is in \$1000s.
\item \ix{deviance}: residual deviance $<$ null deviance $\Rightarrow$ predictors help; lower \ix{AIC} is better.
\end{itemize}
\end{topic}

\begin{trap}[Perfect separation makes logistic regression unstable]
\ixh{perfect separation}If $Y=0$ for $X\le50$ and $Y=1$ for $X>50$ (\textbf{perfectly separable}),
\texttt{glm(y\textasciitilde x, family = "binomial")} warns
``\texttt{algorithm did not converge}'' and ``\texttt{fitted probabilities numerically 0 or 1 occurred}'':
the MLE $\hat\beta_1\to\infty$, so the coefficient estimates are unstable. \textbf{LDA does not suffer
from this problem} \src{S3.45}.
\end{trap}

%% ==================================================================
\subsection{Multiple logistic regression and confounding}
\ixh{multiple logistic regression}\ixh{confounding}

\begin{topic}{The student paradox on the Default data}
\src{S3.27--29}\par\noindent\resizebox{\linewidth}{!}{\begin{tabular}{@{}lrr@{}}
\toprule
 & \texttt{default\textasciitilde student} & \texttt{\textasciitilde balance+income+student}\\ \midrule
Intercept & $-3.5041$ & $-10.8690$\\
balance & --- & $0.0057$ \;($p<.0001$)\\
income (\$1000s) & --- & $0.0030$ \;($p=.7115$)\\
student[Yes] & $\mathbf{0.4049}$ \;($p=.0004$) & $\mathbf{-0.6468}$ \;($p=.0062$)\\ \bottomrule
\end{tabular}}
\begin{itemize}
\item \textbf{Contradicting signs!} Simple model: a student is \textbf{riskier} than a non-student
  \emph{if no other information} about the student is available. Multiple model: a student is
  \textbf{less risky} than a non-student \emph{with the same credit-card balance} (and income).
\item Why: \textbf{student and balance are correlated} --- students are more likely to carry large balances,
  which are associated with high default rates. Balance \ix{confounding} the student effect.
\item Lesson: pay extra attention when regressing on a \textbf{single predictor} while other relevant predictors
  exist (\ix{confounding}; \ix{Simpson's paradox}: S3.30 shows two groups each with a positive
  $y$--$x$ trend, but the pooled data trend is negative) \src{S3.29--31}.
\end{itemize}
\end{topic}

%% ==================================================================
\subsection{Predicted probabilities and the threshold}
\ixh{predicted probability}\ixh{threshold}

\begin{keybox}[Worked prediction (Default data)]
\ixh{predicted probability}$\widehat{\Pr}(\text{default}=\text{Yes}\mid \text{balance}=1500,$ $\text{income}=\$40\text{k},\ \text{student}=\text{No})$
\[
=\frac{e^{-10.869+0.00574(1500)+0.0030(40)-0.6468(0)}}{1+e^{-10.869+0.00574(1500)+0.0030(40)-0.6468(0)}}
\]
\[
=\frac{e^{-2.14}}{1+e^{-2.14}}=0.105
\]
Same person but student $=$ Yes: $z=-2.14-0.6468=-2.79\Rightarrow\hat p=0.058$ \src{S3.32}.
{\footnotesize Recipe: (1) $z=\hat\beta_0+\sum\hat\beta_jx_j$; (2) $\hat p=e^z/(1+e^z)$. The slide writes
``$0.0030\times40000$'', but with income in \$1000s you must plug in 40 (40000 would give $\hat p\approx1$);
fully rounded $\hat\beta$'s give $\approx0.10$.}
\end{keybox}

\begin{topic}{From predicted probability to predicted class}
\src{S3.33}
\begin{itemize}
\item Classify by \textbf{thresholding} the \ix{predicted probability}: predict default $=$ Yes if
  $\hat p>0.5$ (the usual \ix{threshold}; $\hat p>0.5\iff$ logit $>0$).
\item A \textbf{conservative} credit company might predict default $=$ Yes whenever $\hat p>0.1$
  $\Rightarrow$ more individuals flagged: more \ix{true positive}s \emph{and} more \ix{false positive}s.
\item Raising the threshold $\Rightarrow$ fewer positive classifications \src{PE Q15}: $\hat p=.31,.48,.63,.81$ for
  borrowers A--D; threshold $0.5\Rightarrow$ C, D; threshold $0.7\Rightarrow$ only D.
\item Changing the threshold changes the \ix{confusion matrix}, \ix{sensitivity}, \ix{specificity} --- but
  \textbf{not} the \ix{ROC curve} or \ix{AUC} (PE Q18). Details: see the model-evaluation section.
\end{itemize}
\end{topic}

\begin{trap}[PE Q14: what a predicted probability means]
Pneumonia model gives one patient $\hat p=0.78$. Correct: \textbf{(c)} ``the fitted model assigns this patient
a probability of 0.78 of belonging to the pneumonia class.'' NOT ``the model is 78\% accurate for this patient'',
NOT ``definitely has pneumonia'', NOT ``78\% of the symptoms indicate pneumonia'' \src{PE Q14}.
\end{trap}

\begin{topic}{Logistic regression in tidymodels (Lab 3)}
\src{Lab 3}\quad Data: \ixc{Smarket} (outcome \texttt{Direction}: Down/Up; predictors
\texttt{Lag1}--\texttt{Lag5}, \texttt{Volume}) and \ixc{OJ} (outcome \texttt{Purchase}: CH/MM).
\begin{rcode}
library(tidymodels); library(discrim); library(klaR)  (*@\ixch{tidymodels}\ixch{discrim}\ixch{klaR}@*)
set.seed(3435)
oj_split <- initial_split(OJ, prop = 0.70, strata = Purchase)  (*@\ixch{initial\_split()}\ixch{strata}@*)
oj_train <- training(oj_split); oj_test <- testing(oj_split)
oj_recipe <- recipe(Purchase ~ PriceCH + PriceMM + DiscCH + DiscMM +
                      LoyalCH + PctDiscMM + PctDiscCH, data = oj_train)
log_reg <- logistic_reg() %>%                     (*@\ixch{logistic\_reg()}@*)
  set_engine("glm") %>%                           (*@\ixch{set\_engine()}\ixch{glm}@*)
  set_mode("classification")                      (*@\ixch{set\_mode()}@*)
ojlog_wkflow <- workflow() %>%                    (*@\ixch{workflow()}@*)
  add_model(log_reg) %>% add_recipe(oj_recipe)    (*@\ixch{add\_model()}\ixch{add\_recipe()}@*)
ojlog_fit <- fit(ojlog_wkflow, oj_train)          (*@\ixch{fit()}@*)
ojlog_fit %>% tidy()                              (*@\ixch{tidy()}@*)
predict(ojlog_fit, new_data = oj_train, type = "prob")  (*@\ixch{predict()}\ixch{type = "prob"}@*)
augment(ojlog_fit, new_data = oj_train) %>%       (*@\ixch{augment()}@*)
  conf_mat(truth = Purchase, estimate = .pred_class)  (*@\ixch{conf\_mat()}\ixch{.pred\_class}@*)
\end{rcode}
\begin{routput}
  term        estimate std.error statistic  p.value
1 (Intercept)     4.47     2.21      2.02  4.35e- 2
6 LoyalCH        -6.07     0.450   -13.5   1.48e-41
...
# predict(type = "prob"):   .pred_CH .pred_MM
                         1     0.606    0.394
                         2     0.898    0.102
\end{routput}
\begin{itemize}
\item \ixc{predict()} default (\texttt{type = "class"}) returns \ixc{.pred\_class} (threshold 0.5);
  \ixc{type = "prob"} returns one column per level, \texttt{.pred\_<level>} (rows sum to 1).
  \ixc{augment()} attaches \emph{both} to the data. \ixc{tidy()} gives the coefficient table.
\item Smarket: training probabilities all $\approx0.4$--$0.6$ $\Rightarrow$ the model cannot discriminate Up/Down
  (training accuracy 0.525); OJ: probabilities near 0/1, training accuracy 0.833.
\end{itemize}
\end{topic}

\begin{trap}[Which level is ``Y = 1''? (factor levels)]
\ixh{factor level}\ixh{event level}\ixc{glm} (and \ixc{logistic\_reg()} with \texttt{"glm"}) models the probability of the
\textbf{second} factor level; the first is the baseline. OJ levels are (CH, MM), so the coefficients are log odds
of buying \textbf{MM}: \texttt{LoyalCH} $=-6.07$ means more Citrus-Hill loyalty $\Rightarrow$ lower odds of MM.
\textbf{yardstick} metrics, in contrast, treat the \textbf{first} level as the event
(\texttt{event\_level = "first"}). HW3 therefore reorders \texttt{survived} so ``Yes'' is first
(\texttt{factor(survived, levels = c("Yes","No"))}) and uses \ixc{.pred\_Yes} in \ixc{roc\_auc()} ---
but then the glm coefficients describe the log odds of \emph{No} (flip the sign to talk about survival).
\end{trap}

%% ==================================================================
\subsection{Discriminant analysis: Bayes' theorem and the idea}
\ixh{discriminant analysis}\ixh{Bayes' theorem}

\begin{topic}{Why bother with LDA?}
\src{S3.45--46}\quad\textbf{(1)} \ix{logistic regression} coefficient estimates can be \textbf{unstable}
(e.g.\ well/perfectly separated classes); \textbf{(2)} $n$ is \textbf{small} and the distribution of $X$ is
close to \textbf{normal} (in each class); \textbf{(3)} \ix{multi-class classification} ($K>2$) problems.
\end{topic}

\begin{keybox}[Bayes' theorem for classification]
The \ix{Bayes classifier} $\operatorname{argmax}_{k=1,\dots,K}\Pr(Y=k\mid X_1,\dots,X_p)$ is the ``best'' we can do
\src{S3.48}. Logistic regression models $\Pr(Y=1\mid X)$ directly; \ix{LDA} instead uses \ix{Bayes' theorem}
and makes assumptions on $X_1,\dots,X_p\mid Y$ \src{S3.48--49}:
\begin{itemize}
\item \ix{prior probability} $\pi_k=\Pr(Y=k)$;\quad
  \ix{density function} of $X$ in class $k$: $f_k(x)=\Pr(X=x\mid Y=k)$.
\end{itemize}
\[
\underbrace{\Pr(Y=k\mid X=x)}_{\text{\ix{posterior probability}}}
=\frac{\Pr(Y=k,\,X=x)}{\Pr(X=x)}
=\frac{\pi_k\,f_k(x)}{\sum_{j=1}^{K}\pi_j\,f_j(x)}
\]
The denominator is the same for every $k$ $\Rightarrow$ classify to the $k$ with the largest $\pi_kf_k(x)$.\\
\textbf{Discriminant analysis in one sentence:} assume $f_k$ is the \textbf{normal} density, i.e.\
$(X_1,\dots,X_p)\mid Y$ is normally distributed \src{S3.68}.\\
{\footnotesize Example (mine): $\pi_1=0.3,\ \pi_2=0.7$, $f_1(x_0)=0.4,\ f_2(x_0)=0.1$ $\Rightarrow$
$\Pr(Y=1\mid x_0)=\frac{0.12}{0.12+0.07}=0.63$ $\Rightarrow$ class 1, despite the smaller prior.}
\end{keybox}

%% ==================================================================
\subsection{LDA with one predictor ($p=1$)}
\ixh{linear discriminant analysis (LDA)}

\begin{keybox}[LDA, p = 1: model, discriminant function, estimates]
\textbf{Assume} $X\mid Y=k\sim N(\mu_k,\sigma^2)$ --- \textbf{different mean, same variance} \src{S3.69}:
\[ f_k(x)=\frac{1}{\sqrt{2\pi}\,\sigma}\exp\Bigl(-\frac{1}{2\sigma^2}(x-\mu_k)^2\Bigr), \]
\[ \Pr(Y=k\mid X=x)\propto\pi_k\exp\Bigl(-\frac{1}{2\sigma^2}(x-\mu_k)^2\Bigr) \]
($\tfrac{1}{\sqrt{2\pi}\sigma}$ and the denominator are the same for all $k$ and cancel.) Taking logs and dropping
terms that do not depend on $k$, assign $x$ to the class whose \ix{discriminant function}
\[ \delta_k(x)=x\cdot\frac{\mu_k}{\sigma^2}-\frac{\mu_k^2}{2\sigma^2}+\log(\pi_k)\quad\text{is largest} \]
{\footnotesize (S3.70--76 print the middle term as $\mu_k^2/\sigma^2$; the algebra gives $2\sigma^2$, as in ISLR.)}\\
$\mu_1,\dots,\mu_K,\pi_1,\dots,\pi_K,\sigma^2$ are unknown $\Rightarrow$ \textbf{estimate them} \src{S3.71, S3.76}:
\[ \hat\pi_k=\frac{n_k}{n},\qquad \hat\mu_k=\frac{1}{n_k}\sum_{i:y_i=k}x_i,\qquad
\hat\sigma^2=\frac{1}{n-K}\sum_{k=1}^{K}\sum_{i:y_i=k}(x_i-\hat\mu_k)^2 \]
($n_k=$ \# training obs.\ in class $k$; $\hat\sigma^2$ is the \ix{pooled variance}.) The LDA classifier plugs these into
$\hat\delta_k(x)=x\,\hat\mu_k/\hat\sigma^2-\hat\mu_k^2/(2\hat\sigma^2)+\log\hat\pi_k$, which is
\textbf{linear in $x$} --- hence ``linear'' discriminant analysis.
\end{keybox}

\begin{topic}{Decision boundary ($p=1$, $K=2$)}
\ixh{decision boundary}\src{S3.77}
\begin{itemize}
\item \ix{decision boundary}: $\delta_1(x)=\delta_2(x)$. With equal priors $\pi_1=\pi_2$:
  $x^*=\frac{\mu_1+\mu_2}{2}$ (midpoint of the means). Unequal priors (ISLR):
  $x^*=\frac{\mu_1+\mu_2}{2}+\frac{\sigma^2\log(\pi_2/\pi_1)}{\mu_1-\mu_2}$ --- the boundary moves toward the
  \emph{less} likely class, enlarging the region of the more likely one.
\item S3.77: left, densities $N(\mu_1,\sigma^2)$ (green) and $N(\mu_2,\sigma^2)$ (purple) cross at 0 $=$
  \ix{Bayes decision boundary} (dashed); right, histograms of training samples with the \textbf{LDA decision
  boundary} (solid, slightly left of 0) --- LDA estimates the Bayes boundary from data.
\item Intuition (S3.52--75, gene-transcription/drug example): LDA builds a \textbf{new axis} that
  \textbf{maximises} $\dfrac{d^2}{s_1^2+s_2^2}$ --- squared distance between class means over the within-class
  ``scatter'' --- rather than ignoring a predictor (projecting on Gene X alone loses separation). For 3 classes:
  $\frac{d_1^2+d_2^2+d_3^2}{s_1^2+s_2^2+s_3^2}$ with $d_k$ $=$ distance from class centroid to the overall centroid;
  data are projected onto LD1, LD2 (S3.82--85).
\end{itemize}
\end{topic}

%% ==================================================================
\subsection{LDA with several predictors ($p>1$)}
\ixh{linear discriminant analysis (LDA)}

\begin{keybox}[LDA, p > 1]
\textbf{Assume} $(X_1,\dots,X_p)\mid Y=k\sim N(\mu_k,\Sigma)$ --- \textbf{different mean vector, same
\ix{covariance matrix}} $\Sigma$ ($p\times p$, shared by all classes) \src{S3.78}
(\ix{multivariate normal distribution}):
\[ f_k(\mathbf x)=\frac{1}{(2\pi)^{p/2}|\Sigma|^{1/2}}
\exp\Bigl(-\tfrac12(\mathbf x-\mu_k)^T\Sigma^{-1}(\mathbf x-\mu_k)\Bigr) \]
``With a lot of algebra'', the Bayes classifier assigns $\mathbf x$ to the class with the largest
\[ \delta_k(\mathbf x)=\mathbf x^T\Sigma^{-1}\mu_k-\tfrac12\mu_k^T\Sigma^{-1}\mu_k+\log\pi_k . \]
LDA plugs in $\hat\mu_k$, $\hat\Sigma$ (pooled), $\hat\pi_k$: $\hat\delta_k(\mathbf x)$ is \textbf{linear in $\mathbf x$}
(only the term $\mathbf x^T\hat\Sigma^{-1}\hat\mu_k$ involves $\mathbf x$) \src{S3.79}.
Posterior probabilities: $\widehat{\Pr}(Y=k\mid\mathbf x)=e^{\hat\delta_k(\mathbf x)}\big/\sum_j e^{\hat\delta_j(\mathbf x)}$.\\
\textbf{Boundary} \src{S3.81}: Bayes $=\{\mathbf x:\delta_k(\mathbf x)=\delta_j(\mathbf x)\}$, LDA
$=\{\mathbf x:\hat\delta_k(\mathbf x)=\hat\delta_j(\mathbf x)\}$, $k\neq j$ $\Rightarrow$ \textbf{linear
boundary, i.e.\ a hyperplane} (a line when $p=2$; with $K=3$, three line pieces meeting at a point). In S3.81 the
LDA (solid) and Bayes (dashed) boundaries nearly coincide.
\end{keybox}

%% ==================================================================
\subsection{Quadratic discriminant analysis (QDA)}
\ixh{quadratic discriminant analysis (QDA)}

\begin{keybox}[QDA]
Key in discriminant analysis: $(X_1,\dots,X_p)\mid Y$ is normal \src{S3.86}.
\textbf{LDA:} $N(\mu_k,\Sigma)$ --- different mean, \textbf{same} covariance.
\textbf{\ix{QDA}:} $N(\mu_k,\Sigma_k)$ --- different mean, \textbf{different} covariance (each class its own $\Sigma_k$).
Assign $\mathbf x$ to the class with the largest \src{S3.87}
\begin{align*}
\delta_k(\mathbf x)=&-\tfrac12\mathbf x^T\Sigma_k^{-1}\mathbf x+\mathbf x^T\Sigma_k^{-1}\mu_k\\
&-\tfrac12\mu_k^T\Sigma_k^{-1}\mu_k-\tfrac12\log|\Sigma_k|+\log\pi_k
\end{align*}
(QDA: hats on $\mu_k,\Sigma_k,\pi_k$.) The term $-\tfrac12\mathbf x^T\hat\Sigma_k^{-1}\mathbf x$ no longer cancels
across classes $\Rightarrow$ \textbf{quadratic in $\mathbf x$} $\Rightarrow$ curved (quadratic)
\ix{decision boundary}.\\
\textbf{Parameter count} (covariances): LDA one $\Sigma$: $p(p+1)/2$; QDA: $K\cdot p(p+1)/2$
(e.g.\ $p=50$: 1275 vs.\ $1275K$) \src{S3.91}.
\end{keybox}

\begin{topic}{S3.88: LDA vs QDA boundaries (purple dashed $=$ Bayes, black dotted $=$ LDA, green $=$ QDA)}
\begin{itemize}
\item $\Sigma_1=\Sigma_2$: the Bayes boundary is \textbf{linear}; LDA is close to it, QDA's curve deviates
  (extra variance, no bias gain).
\item $\Sigma_1\neq\Sigma_2$: the Bayes boundary is \textbf{curved}; QDA tracks it, LDA's straight line cannot (bias).
\end{itemize}
\end{topic}

%% ==================================================================
\subsection{Comparing classifiers}
\ixh{model comparison}

\begin{keybox}[KNN vs logistic vs LDA vs QDA vs naive Bayes]
\ixh{parametric model}\ixh{non-parametric model}\ixh{decision boundary}
{\footnotesize
\begin{tabularx}{\linewidth}{@{}l c l X@{}}
\toprule
Method & Param.? & Boundary & Models / assumes\\ \midrule
\ix{KNN} & no & nonlinear (any shape) & majority vote of the $k$ nearest training points; no assumption on $f$\\
Logistic & yes & linear & $\Pr(Y=1\mid X)$ directly via the logit (\ix{discriminative model}); fit by MLE\\
LDA & yes & linear & $X\mid Y=k\sim N(\mu_k,\Sigma)$, common $\Sigma$ (\ix{generative model})\\
QDA & yes & quadratic & $X\mid Y=k\sim N(\mu_k,\Sigma_k)$, class-specific $\Sigma_k$\\
Naive Bayes & yes & quadratic, no cross terms & predictors independent within each class\\ \bottomrule
\end{tabularx}}
\textbf{Model \ix{flexibility}} \src{S3.89}: KNN $>$ QDA $>$ logistic regression $=$ LDA
(choose via the \ix{bias-variance tradeoff}); \ix{interpretability} runs the other way. KNN is
\ix{non-parametric model}; logistic, LDA, QDA are \ix{parametric model}s.
\end{keybox}

\begin{topic}{Logistic regression vs LDA}
\src{S3.90}
\begin{itemize}
\item \textbf{Both produce linear decision boundaries}; the coefficient-estimation procedures differ
  (logistic: \ix{maximum likelihood}; LDA: plug in estimated normal means/variance/priors).
\item LDA assumes a \textbf{normal} distribution of the predictors given the response class.
\item Logistic regression learns $\Pr(Y=k\mid X)$ \textbf{directly}: \ix{discriminative model} (discriminative learning).
  LDA learns $\Pr(Y=k\mid X)$ by modelling $\Pr(Y=k,X)=\pi_kf_k(x)$: \ix{generative model} (generative learning).
\item Rule of thumb: normality roughly holds / small $n$ / well-separated classes / $K>2$ $\Rightarrow$ LDA;
  normality doubtful $\Rightarrow$ logistic regression is more robust.
\end{itemize}
\end{topic}

\begin{topic}{LDA vs QDA (bias--variance)}
\src{S3.91}
\begin{itemize}
\item LDA is a \textbf{special case} of QDA (all $\Sigma_k=\Sigma$). QDA is \textbf{more flexible}, but has
  \textbf{more parameters} to estimate (the $\Sigma_k$'s) $\Rightarrow$ lower \ix{bias}, higher \ix{variance}.
\item Prefer \textbf{LDA}: few training observations (small $n$) relative to $p$, or class covariances similar
  --- reducing variance is crucial. Prefer \textbf{QDA}: large $n$ (variance less of a concern) or the
  common-covariance assumption is clearly untenable.
\item In practice we often do not know which assumptions are reasonable --- let the data decide via
  \ix{cross-validation} \src{S3.93}.
\end{itemize}
\end{topic}

\begin{qbox}{PE Q20 (LDA vs QDA)}
50{,}000 observations, two predictors, clear evidence that the covariance structures differ substantially
across outcome levels. Which has greater potential to capture the true decision boundary?
(a) LDA (b) QDA (c) equal (d) neither (both parametric).
\ans \textbf{(b) QDA.} LDA assumes a common covariance, so it would not do well; QDA estimates more
parameters, but $p=2$ is small and $n$ is huge, so the extra variance is negligible (parametric is fine)
\src{PE Q20}.
\end{qbox}

\begin{trap}[LDA/QDA traps]
\begin{itemize}
\item \textbf{QDA only helps if the class covariances really differ}; if $\Sigma_1\approx\Sigma_2$ it just adds
  variance (more parameters, no bias reduction) and LDA usually wins on test data.
\item ``Linear'' in LDA refers to the \textbf{decision boundary} (linear $\hat\delta_k$), not to linear regression.
  LDA and logistic regression share the boundary shape but are estimated differently.
\item LDA/QDA assume \textbf{normal} predictors within each class --- categorical/dummy predictors violate this
  (they still run, but the assumption is wrong).
\item Higher \ix{training error} rate is not the criterion: the most flexible model (QDA/KNN) tends to win on
  \emph{training} data and can lose on \ix{testing data} (\ix{overfitting}).
\end{itemize}
\end{trap}

%% ==================================================================
\subsection{Naive Bayes (Lab 3)}
\ixh{naive Bayes}

\begin{topic}{Naive Bayes}
\begin{itemize}
\item Same Bayes'-theorem idea, $\Pr(Y=k\mid x)\propto\pi_kf_k(x)$, but instead of a multivariate normal,
  assume the $p$ predictors are \textbf{independent within each class}:
  $f_k(x)=f_{k1}(x_1)\times f_{k2}(x_2)\times\dots\times f_{kp}(x_p)$.
\item With \ixc{usekernel} \texttt{= FALSE} each $f_{kj}$ is a \textbf{normal (Gaussian)} density
  (\texttt{TRUE} $\Rightarrow$ a kernel density estimate) \src{Lab 3}. Gaussian naive Bayes $\approx$ QDA with
  \emph{diagonal} $\Sigma_k$ (no within-class correlations).
\item ``Naive'' $=$ the independence assumption (usually false); it slashes the number of parameters
  $\Rightarrow$ higher bias, much lower variance; useful when $p$ is large relative to $n$.
\end{itemize}
\end{topic}

\begin{topic}{LDA / QDA / naive Bayes in tidymodels (Lab 3)}
\ixch{discrim}Needs \ixc{library(discrim)} (and \ixc{klaR} for naive Bayes). The recipe/workflow is reused; only the
model spec changes:
\begin{rcode}
lda_mod <- discrim_linear() %>%                     (*@\ixch{discrim\_linear()}@*)
  set_mode("classification") %>% set_engine("MASS") (*@\ixch{MASS}\ixch{set\_engine()}@*)
qda_mod <- discrim_quad() %>%                       (*@\ixch{discrim\_quad()}@*)
  set_mode("classification") %>% set_engine("MASS")
nb_mod <- naive_Bayes() %>%                         (*@\ixch{naive\_Bayes()}@*)
  set_mode("classification") %>% set_engine("klaR") %>%  (*@\ixch{klaR}@*)
  set_args(usekernel = FALSE)                       (*@\ixch{set\_args()}\ixch{usekernel}@*)
smarketqda_wkflow <- workflow() %>% add_model(qda_mod) %>%
  add_recipe(smarket_recipe)
smarketqda_fit <- fit(smarketqda_wkflow, smarket_train)
augment(smarketqda_fit, new_data = smarket_train) %>%
  accuracy(truth = Direction, estimate = .pred_class)    (*@\ixch{accuracy()}@*)
\end{rcode}
\textbf{Spelling:} \ixc{naive\_Bayes()} (capital B), \ixc{discrim\_linear()}, \ixc{discrim\_quad()};
engines \texttt{"MASS"} (LDA, QDA), \texttt{"klaR"} (naive Bayes).

\textbf{Lab 3 training accuracies} (picked by training accuracy only because CV had not been taught):\\[1pt]
{\footnotesize
\begin{tabular}{@{}lcc@{}}
\toprule
Model & \ixc{Smarket} & \ixc{OJ}\\ \midrule
Logistic regression & 0.525 & \textbf{0.833}\\
LDA & 0.524 & 0.828\\
QDA & \textbf{0.548} & 0.797\\
Naive Bayes & 0.529 & 0.778\\ \bottomrule
\end{tabular}}\\[1pt]
\textbf{Test set:} Smarket QDA accuracy \textbf{0.487} (sensitivity 0.254, specificity 0.703) --- worse than a
coin flip despite the best training accuracy; OJ logistic accuracy 0.832 (sens.\ 0.908, spec.\ 0.714).
Lesson: choosing a model by \ix{training error} is \textbf{not} a good idea in general; use
\ix{cross-validation}. Test metrics via \ixc{metric\_set()}\texttt{(accuracy, sensitivity, specificity)}
and \ixc{roc\_curve()}\texttt{(Direction, .pred\_Down)} (see model evaluation).
\end{topic}

%% ==================================================================
\subsection{HW3: four classifiers on the Titanic data}
\ixch{titanic}

\begin{topic}{HW3 pipeline (Q1, Q4--Q10)}
\src{HW3}\quad \ixc{titanic} (891 passengers; \texttt{survived} Yes $\approx38\%$; \texttt{age} has 177 NAs).
Make \texttt{survived} and \texttt{pclass} factors (``Yes'' first), split with \ixc{strata} \texttt{= survived}
(\ix{stratified sampling} keeps the $\approx$38/62 class ratio in both sets --- \ix{class imbalance}).
\begin{rcode}
titanic <- read_csv("data/titanic.csv") %>%                (*@\ixch{read\_csv()}@*)
  mutate(survived = factor(survived, levels = c("Yes", "No")),  (*@\ixch{factor()}@*)
         pclass = factor(pclass))
set.seed(131)
t_split <- initial_split(titanic, prop = 0.7, strata = survived)
t_train <- training(t_split); t_test <- testing(t_split)
t_recipe <- recipe(survived ~ pclass + sex + age + sib_sp + parch + fare,
                   data = t_train) %>%
  step_impute_linear(age, impute_with = imp_vars(sib_sp, parch, fare)) %>%  (*@\ixch{step\_impute\_linear()}\ixch{imp\_vars()}@*)
  step_dummy(all_nominal_predictors()) %>%                 (*@\ixch{step\_dummy()}\ixch{all\_nominal\_predictors()}@*)
  step_interact(terms = ~ starts_with("sex"):fare + age:fare)  (*@\ixch{step\_interact()}\ixch{starts\_with()}@*)
log_spec <- logistic_reg() %>% set_engine("glm") %>% set_mode("classification")
lda_spec <- discrim_linear() %>% set_engine("MASS") %>% set_mode("classification")
qda_spec <- discrim_quad() %>% set_engine("MASS") %>% set_mode("classification")
knn_spec <- nearest_neighbor(neighbors = 10) %>%           (*@\ixch{nearest\_neighbor()}\ixch{neighbors}@*)
  set_engine("kknn") %>% set_mode("classification")        (*@\ixch{kknn}@*)
log_fit <- workflow() %>% add_model(log_spec) %>% add_recipe(t_recipe) %>%
  fit(t_train)                                  # repeat for lda/qda/knn
predict(log_fit, new_data = t_train, type = "prob") %>%    (*@\ixch{type = "prob"}@*)
  bind_cols(t_train %>% select(survived)) %>%             (*@\ixch{bind\_cols()}@*)
  roc_auc(truth = survived, .pred_Yes)                     (*@\ixch{roc\_auc()}\ixch{.pred\_Yes}@*)
augment(log_fit, new_data = t_test) %>%
  conf_mat(truth = survived, estimate = .pred_class) %>% autoplot(type = "heatmap")  (*@\ixch{autoplot()}@*)
augment(log_fit, new_data = t_test) %>% roc_curve(survived, .pred_Yes) %>% autoplot()  (*@\ixch{roc\_curve()}@*)
\end{rcode}
\begin{itemize}
\item Order of steps: \textbf{impute} $\to$ \textbf{dummy} $\to$ \textbf{interact}. After \ixc{step\_dummy()}
  \texttt{sex} becomes \texttt{sex\_male}, so the \ix{interaction term} must use \ixc{starts\_with()}\texttt{("sex")}.
  (\ixc{imp\_vars()} lists the predictors of the linear \ix{imputation} model; use complete columns.)
\item \ix{AUC} needs \textbf{probabilities} (\texttt{type = "prob"}, column \texttt{.pred\_Yes}), not
  \texttt{.pred\_class}. Compare \textbf{training vs testing} AUC: a big drop (typical for KNN with small $k$,
  and for the more flexible QDA) signals \ix{overfitting}.
\item HW3 Q11--12 (231): see the logit-inverse box above.
\end{itemize}
\end{topic}
